
\section{Introduction}


\begin{frame}
\frametitle{Infrastructures \textit{cloud}}

Les données sont massives quand elles dépassent les capacités d'une
seule machine.

\begin{redbullet}
  \item en mémoire RAM: quelques dizaines de GOs;
  \item en mémoire persistante: quelques TOs.
\end{redbullet}

\vfill

Qui dit \alert{données massives} dit donc \alert{système distribué}. Point communs (le seul?) des 
systèmes dits NoSQL.

\vfill

De fait: l'infrastructure est constituée de grappes de serveurs, avec
\begin{redbullet}
  \item des machines à bas coût (\textit{commodity servers}),
  \item la possibilité d'en ajouter/retirer en quelques minutes (\alert{élasticité}),
  \item des environnements logiciels spécialisés.
\end{redbullet}

Nous appellerons l'ensemble un \textit{cloud} (un peu réducteur).


\end{frame}


\section{Cloud et données massives}

\begin{frame}{Vision générale}

\figSlideWithSize{cloud}{12cm}

\end{frame}

\begin{frame}{Serveurs et baies}

Dans un centre de données, les serveurs sont empilés dans des baies.
\vfill

\figSlideWithSize{rack-and-server}{11cm}

\end{frame}


\begin{frame}
\frametitle{Infrastructures \textit{cloud}}



% A two-columns presentation
\begin{tabular}{l|p{6cm}}
   \parbox{5cm}{

Une baie contient environ 40 serveurs.

\smallskip

Un centre de données contient quelques centaines de baies.

\smallskip

Ordre de grandeur: quelques milliers de serveurs par centre.

\smallskip

Combien des centres chez Google, Facebook, Amazon ...? Des millions
de serveurs.


   }  % Inclusion of an XML file
  & 
  \parbox{6cm}{\figSlideWithSize{google-datacenter}{6cm} }
\end{tabular}
\vfill

\end{frame}


\begin{frame}
\frametitle{La virtualisation}

Prenez une machine avec 64GOs de mémoire, 12 disques de 3 TO chacun.

\medskip

Découpez-là en sous-machines ``virtuelles''.

\begin{redbullet}
  \item 18 serveurs virtuels avec 4 GO de RAM chacun (2 TO de disques)\,;
  \item 8 serveurs virtuels avec 8 GO de RAM chacun (3 TO de disques)\,;

    \item etc...
\end{redbullet}

Permet d'attribuer des serveurs à la demande, optimise l'utilisation des machines.

\medskip

Pour le client\,: 0 coût installation, maintenance, amortissement.

\medskip

Cf. Amazon EC2, OVH, Gandi (français), etc.

\end{frame}


\begin{frame}
\frametitle{Les pannes}

Choix essentiel: \alert{on augmente la puissance du système par ajout de machines
à bas coût (scalabilité horizontale)}, et pas par l'augmentation des capacités
d'une machine (scalabilité verticale).

\vfill

Conséquence essentielle: des pannes tout le temps!

\begin{redbullet}
  \item panne de disque;
  \item redémarrage serveur (panne logicielle/matérielle)
  \item coupure réseau.
\end{redbullet}

Méthode générale: réplication des données; redondance des services.

\vfill

Les systèmes (NoSQL) disposent d'une reprise sur panne (\textit{failover}) \alert{automatique}.

\end{frame}

\section{Systèmes distribués}



\begin{frame}
\frametitle{Systèmes distribués}

Système distribué = ensemble de composants logiciels, ou \alert{nœuds}, répartis sur une \alert{grappe de serveurs},
coopérant pour une tâche précise.

\vfill

Ces nœuds communiquent par \alert{envoi de messages} (aucun partage de mémoire). 


\vfill

Chaque nœud est en écoute sur un port réseau (ex., le 80 pour le Web, 27017 pour Mongodb).

\vfill

On peut avoir un système distribué (mais pas scalable) sur une seule machine! 

\vfill
 En général,
un nœud par machine.

\end{frame}


\begin{frame}
\frametitle{Le client, les maîtres et les esclaves}

Deux organisations principales. La première (plus courante) est
dite \alert{maître-esclave}.

\vfill


\figSlideWithSize{client-master-slave}{11cm}

\vfill

Rôle essentiel du maître (mais aussi \textit{single point of failure}).


\end{frame}


\begin{frame}
\frametitle{Autre topologie}

C'est la démocratie! Il n'y a que des maîtres (\textit{master-master}, ou aussi multi-nœuds).

\vfill


\figSlideWithSize{client-master}{11cm}

\vfill

Plus robuste en théorie, mais soulève des problèmes compliqués (donc, moins courant).

\end{frame}


\begin{frame}{Gestion des pannes, ou \textit{failover}}

Un système distribué peut s'appuyer  des centaines, des milliers de serveurs: \alert{il y a des pannes
tout le temps}.

\vfill

Caractéristique d'un bon système: tolérant aux pannes, assure une disponibilité 24/7.

\begin{redbullet}
  \item surveillance automatisée de tous les nœuds participants (\textit{heartbeat}, messages toutes les $n$ secondes).
  \item \alert{méthode de reprise sur panne} (\textit{failover}) basée sur la redondance/réplication.
\end{redbullet}

\vfill

Un cas particulièrement épineux: le \alert{partitionnement réseau}.

\end{frame}

\section{Le NoSQL}


\begin{frame}{Le NoSQL, qu'est-ce que c'est}

\begin{defblock}{Systèmes NoSQL (Ma définition)}
Un système NoSQL est un système distribué dont la tâche est la gestion de données persistantes
dans un environnement \textit{cloud}. Il fournit les fonctionnalités suivantes:

\begin{redbullet}
  \item Recherche et mise à jour de données 
  \item Adaptation automatique aux ressources matérielles: équilibrage, élasticité.
  \item Performances ``scalables'' (voir plus loin).
  \item Reprise sur panne automatisée.
\end{redbullet}
\end{defblock}

Beaucoup de modèles différents (``documents'', clé-valeur, graphes ...). Deux
types principaux:

\begin{redbullet}
  \item Systèmes \alert{temps réel}: accès en millisecondes aux unités d'information.
  \item Systèmes \alert{analytiques}: traitements appliqués à tout ou partie d'une collection.
\end{redbullet}

\end{frame}


\begin{frame}{Petit historique: systèmes analytiques}

Quelques publications de Google: \textit{Google File System} (2003),
MapReduce (2004), BigTable (2006).

\vfill

Clonés par des systèmes Open Source dans l'environnement Hadoop: HDFS, MapReduce, HBase.

\vfill

Autres successeurs: Cassandra, maintenant Spark, Flink, etc.

\vfill
Applications: construction d'index, analyses de données massives.
 
\end{frame}

\begin{frame}{Petit historique: systèmes temps réels}

Publication majeure d'Amazon: le système Dynamo (2007).

\vfill

Application: gestion robuste, distribuée, scalable de centaines de millions de comptes
client, produits, transactions.

\vfill

Cloné par le système Voldemort; principes repris par Mongodb, CouchDB, Riak, beaucoup d'autres

\vfill

NB: précédé par beaucoup de recherche et de systèmes distribués bien sûr.

\end{frame}

\begin{frame}{Ce qu'il faut retenir}

Environnement dit \textit{cloud} = grappes de serveurs à bas coût, stockés dans des centres de données.

\vfill

Très ``élastique'' car facile d'ajouter/supprimer une ressource = \alert{scalabilité horizontale}.

\vfill

Très sujet aux \alert{pannes}, de disque, de réseau, de logiciels.

\vfill

Système distribué (NoSQL): système exploitant au mieux les ressources du \textit{cloud}, tolérant aux pannes.

 
\end{frame}

\end{document}


